We now treat some special cases. If \(N=\Gamma_0(R)\), write
\[
 \cR_N=\operatorname{ad}(\cR).
\]
If \(N=\VR(R)\cap M\), write
\[
 \cR_N=\operatorname{ss}(\cR).
\]
Thus there is a diagram
\[
 \operatorname{ad}(\cR)\longrightarrow
 \operatorname{ss}(\cR)\longrightarrow\cR.
\]
Put
\[
 \operatorname{der}(\cR)=\operatorname{ss}(\cR^\ast)^\ast,
 \qquad
 \operatorname{sc}(\cR)=\operatorname{ad}(\cR^\ast)^\ast.
\]
By duality there is a diagram
\[
 \cR\longrightarrow
 \operatorname{der}(\cR)\longrightarrow
 \operatorname{sc}(\cR).
\]

\begin{proposition}\label{sga3:XXI:6.5.5}
\begin{enumerate}
\item[(i)] In the first row of the diagram
\begin{center}
\begin{tikzcd}[column sep=large,row sep=large]
\operatorname{ad}(\cR) \arrow[r]
  & \operatorname{ss}(\cR) \arrow[rr] \arrow[dr]
  && \operatorname{der}(\cR) \arrow[r]
  & \operatorname{sc}(\cR) \\
&& \cR \arrow[ur] &&
\end{tikzcd}
\end{center}
the four root data are semisimple and the three morphisms are
isogenies.
\item[(ii)] The datum \(\operatorname{ad}(\cR)\) is adjoint, and
\(\cR\) is adjoint if and only if
\(\operatorname{ad}(\cR)\to\cR\) is an isomorphism.
\item[(iii)] The datum \(\operatorname{sc}(\cR)\) is simply connected,
and \(\cR\) is simply connected if and only if
\(\cR\to\operatorname{sc}(\cR)\) is an isomorphism.
\item[(iv)] The following conditions are equivalent:
\begin{enumerate}
\item[(a)] \(\cR\) is semisimple;
\item[(b)] \(\operatorname{ss}(\cR)\to\cR\) is an isomorphism;
\item[(c)] \(\cR\to\operatorname{der}(\cR)\) is an isomorphism.
\end{enumerate}
\end{enumerate}
\end{proposition}

Consider the morphism
\[
 \operatorname{ss}(\cR)\longrightarrow\operatorname{der}(\cR).
\]
The definitions immediately give the following result.

\begin{lemma}\label{sga3:XXI:6.5.6}
Let
\[
 h:\operatorname{ss}(\cR)\longrightarrow\operatorname{der}(\cR)
\]
be the canonical isogeny. Then
\[
 K(h)\simeq N(\cR).
\]
\end{lemma}

\SGARefTarget{sga3:XXI:6.5.7}
Consider other special cases of induced root data. Put
\[
 N=
 \{x\in M\mid(\alpha^\ast,x)=0\text{ for every }\alpha\in R\}
 \times\Gamma_0(R).
\]
The sum is direct by \ref{sga3:XXI:1.2.5}. It follows that
\(\cR_N\) is identified with
\[
 \operatorname{ad}(\cR)\times\operatorname{corad}(\cR).
\]
One may replace \(\Gamma_0(R)\) by \(\VR(R)\cap M\), and then dualize
both constructions. This gives the following diagram of root data:
\begin{center}
\resizebox{\textwidth}{!}{%
\begin{tikzpicture}[
  >=stealth,
  every node/.style={font=\small},
  x=1cm,y=1cm
]
\node (at) at (0,2.1) {$\operatorname{ad}(\cR)$};
\node (st) at (4,2.1) {$\operatorname{ss}(\cR)$};
\node (dt) at (10,2.1) {$\operatorname{der}(\cR)$};
\node (ct) at (14,2.1) {$\operatorname{sc}(\cR)$};

\node (am) at (0,0) {$\operatorname{ad}(\cR)\times\operatorname{corad}(\cR)$};
\node (sm) at (4,0) {$\operatorname{ss}(\cR)\times\operatorname{corad}(\cR)$};
\node (rm) at (7,0) {$\cR$};
\node (dm) at (10,0) {$\operatorname{der}(\cR)\times\operatorname{rad}(\cR)$};
\node (cm) at (14,0) {$\operatorname{sc}(\cR)\times\operatorname{rad}(\cR)$};

\node (ab) at (0,-2.1) {$\operatorname{ad}(\cR)$};
\node (sb) at (4,-2.1) {$\operatorname{ss}(\cR)$};
\node (db) at (10,-2.1) {$\operatorname{der}(\cR)$};
\node (cb) at (14,-2.1) {$\operatorname{sc}(\cR)$};

\draw[->] (at) -- (st);
\draw[->] (st) -- (dt);
\draw[->] (dt) -- (ct);
\draw[->] (am) -- (sm);
\draw[->] (sm) -- (rm);
\draw[->] (rm) -- (dm);
\draw[->] (dm) -- (cm);
\draw[->] (ab) -- (sb);
\draw[->] (sb) -- (db);
\draw[->] (db) -- (cb);

\draw[->] (at) -- (am);
\draw[->] (am) -- (ab);
\draw[->] (st) -- (sm);
\draw[->] (sm) -- (sb);
\draw[->] (dt) -- (dm);
\draw[->] (dm) -- (db);
\draw[->] (ct) -- (cm);
\draw[->] (cm) -- (cb);
\end{tikzpicture}
}
\end{center}
It is commutative, as is checked at once, and is self-dual in the
evident sense. The horizontal morphisms are isogenies, and the
composites of the vertical arrows are the identity.

\begin{lemma}\label{sga3:XXI:6.5.8}
Let \(h_1\) and \(h_2\) be the canonical isogenies
\[
 \operatorname{ss}(\cR)\times\operatorname{corad}(\cR)
 \xrightarrow{h_1}\cR
 \xrightarrow{h_2}
 \operatorname{der}(\cR)\times\operatorname{rad}(\cR).
\]
Then
\[
 K(h_1)\simeq K(h_2)\simeq N(\cR).
\]
\end{lemma}

\begin{proof}
This is immediate from the definitions.
\end{proof}

\begin{corollary}\label{sga3:XXI:6.5.9}
For a root datum \(\cR\), the following conditions are equivalent:
\begin{enumerate}
\item[(i)] \(N(\cR)=0\), that is,
  \(\operatorname{corad}(\cR)\to\operatorname{rad}(\cR)\) is an
  isomorphism;
\item[(ii)] the canonical map
  \(\operatorname{ss}(\cR)\to\operatorname{der}(\cR)\) is an
  isomorphism;
\item[(iii)] \(h_1\) is an isomorphism;
\item[(iv)] \(h_2\) is an isomorphism;
\item[(v)] \(\cR\) is the product of a semisimple root datum and a
  trivial root datum.
\end{enumerate}
\end{corollary}

\begin{corollary}\label{sga3:XXI:6.5.10}
For every root datum \(\cR\), there are isogenies
\[
 \operatorname{ad}(\cR)\times\cR_0
 \longrightarrow\cR\longrightarrow
 \operatorname{sc}(\cR)\times\cR_0,
\]
where \(\cR_0\) is the trivial root datum of rank
\(\rgred(\cR)-\rgss(\cR)\).
\end{corollary}

\begin{lemma}\label{sga3:XXI:6.5.11}
Let \(\cR\) be a root datum, let \(\Delta\) be a system of simple
roots, and let \(\Delta'\subset\Delta\). Consider the root datum
\[
 \cR_{\Delta'}=(M,M^\ast,R_{\Delta'},R^\ast_{\Delta'})
\]
of \ref{sga3:XXI:3.4.7}.
\begin{enumerate}
\item[(i)] If \(\cR\) is simply connected, then
  \(\operatorname{der}(\cR_{\Delta'})\) is simply connected.
\item[(ii)] If \(\cR\) is adjoint, then
  \(\operatorname{ss}(\cR_{\Delta'})\) is adjoint.
\end{enumerate}
\end{lemma}

\begin{proof}
The two assertions are equivalent by duality. The second reduces to
\[
 M\cap\VR(R_{\Delta'})=\Gamma_0(R_{\Delta'}).
\]
When \(M=\Gamma_0(R)\), both sides are the subgroup of \(M\) generated
by \(\Delta'\).
\end{proof}
